\section{Experiments}
\label{sec:experiments}

\subsection{Dataset}
\label{sec:exp_data}

We evaluate the proposed method on the BCI Competition~IV-2a dataset, a standard
benchmark for multi-class motor-imagery EEG analysis. Its well-structured,
multi-subject, multi-session design makes it well suited to assessing
cross-subject modeling and the generalization of denoising algorithms.

The dataset contains EEG recordings from nine healthy subjects (A01--A09). Each
subject completed two sessions on different days (Session~T and Session~E), used
respectively for training and testing. Each session comprises six runs of 48
trials, giving 288 trials per subject per session. Every trial follows a uniform
protocol: the subject first fixates on a central cross (0--2\,s); a directional
cue (left, right, down, or up, corresponding to imagined movement of the left
hand, right hand, both feet, or tongue) then appears for about 1.25\,s; the
subject performs the corresponding motor imagery until about 6\,s into the
trial, followed by a short rest. Each trial lasts about 8\,s and comprises a
preparation, execution, and recovery phase. Such motor imagery elicits
event-related desynchronization and synchronization (ERD/ERS) over the
sensorimotor cortex~\cite{pfurtscheller1999erd}.

EEG was recorded with 22 channels in a Cz-referenced montage at a 250\,Hz
sampling rate; three additional EOG channels were recorded to support artifact
analysis. All signals were hardware-filtered (0.5--100\,Hz) and acquired with
high-quality equipment to ensure consistency. We use the data of all nine
subjects and adopt the official split without manual resampling or truncation,
keeping Session~T (training) and Session~E (testing) strictly separate to
emulate the independent-test scenario of real deployment. Because the dataset
covers four motor-imagery classes and contains natural artifacts (e.g., eye
blinks and muscle activity), it provides a solid basis for studying EEG
denoising under complex noise.

\subsection{Data Preprocessing}
\label{sec:exp_preprocess}

We apply standardized preprocessing to improve signal quality and ensure
consistency across subjects. For channel selection, we retain all 22 EEG
channels and discard the reference and EOG channels to focus on cortical
activity. We then apply a 1--50\,Hz finite-impulse-response (FIR) band-pass
filter to remove low-frequency baseline drift and high-frequency EMG
interference while preserving the cognitively and motor-relevant bands, followed
by a 50\,Hz notch filter to suppress power-line interference and its harmonics.
Recordings with a sampling rate above 250\,Hz are resampled to 250\,Hz to keep
the input dimension consistent and reduce computational load. We segment the
continuous EEG with a sliding window of 2\,s and a step of 0.5\,s, generating a
sufficient number of training samples while maintaining coverage. Finally, to
remove distribution shifts caused by differing amplitude scales across subjects,
each segment is normalized per channel to zero mean and unit variance, which
stabilizes training and improves cross-subject generalization.

\subsection{Experimental Setup}
\label{sec:exp_setup}

To ensure stable learning under complex noise and across subjects, \saddpm{}
uses a unified, structured set of hyperparameters. We optimize with Adam at an
initial learning rate of $1\times10^{-4}$, combined with a cosine-annealing
schedule for smooth decay and stable late-stage convergence. The batch size is
64, and the diffusion process uses $T=1000$ steps with a linear schedule for the
noise level $\beta_t$, ensuring a gradual transition from the noisy state to the
empirical target representation. Training runs for 100 epochs over the source
Session~T set. For source-domain individual modeling, the source subject is
assigned a 128-dimensional learnable embedding that captures subject-dependent
signal and artifact characteristics; this embedding is optimized jointly with
the backbone and injected at multiple stages of the reverse process. To improve
generalization and avoid overfitting the embeddings to the training data, we
apply $L_2$ regularization with weight $1\times10^{-4}$ to the embedding
vectors. When the evaluation target differs from the source, the target
embedding is not assumed to exist beforehand; it is adapted from unlabeled
target-subject EEG windows as described below. All experiments run on the same
hardware (an NVIDIA~A100 GPU), with identical training configurations across
methods to ensure fair comparison. The overall source-training procedure is
summarized in Algorithm~\ref{alg:training}.

\begin{algorithm}[t]
\caption{\saddpm{} source training.}
\label{alg:training}
\begin{algorithmic}[1]
\Require Source training set of preprocessed EEG segments with subject labels $\{(\mathbf{x}_0,s)\}\sim q(\mathbf{x}_0,s)$; noise schedule $\{\beta_t\}_{t=1}^{T}$
\State Initialize the source embedding(s); if more than one source subject is used, pre-train the subject classifier $\omega$ on source labels
\While{not converged}
  \State Sample a source training segment and its subject label $(\mathbf{x}_0,s)\sim q(\mathbf{x}_0,s)$
  \State Sample a diffusion step $t\sim\mathrm{Uniform}(\{1,\dots,T\})$
  \State Sample noise $\boldsymbol{\epsilon}\sim\mathcal{N}(\mathbf{0},\mathbf{I})$
  \State Form the noisy state $\mathbf{x}_t=\sqrt{\bar{\alpha}_t}\,\mathbf{x}_0+\sqrt{1-\bar{\alpha}_t}\,\boldsymbol{\epsilon}$
  \State Compute
  $\boldsymbol{\epsilon}^{c}_{\theta}(\mathbf{x}_t,t)$,
  $\boldsymbol{\epsilon}^{s}_{\phi}(\mathbf{x}_t,t,\mathbf{e}(s))$,
  $\widehat{\boldsymbol{\epsilon}}_{\theta,\phi}$,
  and the subject logits.
  \State Compute $\mathcal{L}=\lambda_r\mathcal{L}_r+\lambda_o\mathcal{L}_o+\lambda_a\mathcal{L}_a$, with $\lambda_a=0$ in single-source training
  \State Update $\theta,\phi,\omega$ and the source embedding(s) by a gradient-descent step on $\mathcal{L}$
\EndWhile
\end{algorithmic}
\end{algorithm}

To assess cross-subject transfer, we adopt a pairwise source--target protocol
with transductive target-subject adaptation. For each row of the matrix, the
shared \saddpm{} backbone and the source embedding are trained only on the
Session~T data of one source subject. For each target column, before denoising
Session~E, we freeze the backbone and optimize a target embedding from unlabeled
target Session~T EEG windows using Eq.~\eqref{eq:target_embedding_adapt};
no target motor-imagery labels, clean targets, or artifact annotations are used.
In this single-source setting, the source-embedding mean used to initialize the
adapted target embedding reduces to the learned source embedding. The adapted
embedding is then used in the trial-wise partial-noising restoration procedure
of Eqs.~\eqref{eq:partial_noising}--\eqref{eq:trialwise_output}. The
restoration strength $t_0$ and the number of reverse samples $K$ are fixed
across all source--target pairs.

The downstream four-class motor-imagery classifier is trained on denoised source
Session~T trials and evaluated on denoised target Session~E trials; target
labels are used only to compute the final classification accuracy. The resulting
$9\times9$ accuracy matrix has diagonal entries for session-separated
within-subject transfer and off-diagonal entries for cross-subject transfer.
For each target (column) subject, the row labeled ``Mean'' reports the average
accuracy over all source subjects, summarizing how well denoised representations
transfer to that target. This protocol is transductive rather than fully
zero-shot: it evaluates whether \saddpm{} can adapt to a deployment
subject's unlabeled signal distribution while preserving a label-free target setting.

\subsection{Results and Analysis}
\label{sec:exp_results}

To evaluate the effectiveness of \saddpm{} for denoising motor-imagery EEG, we
conduct systematic experiments on BCI Competition~IV-2a and compare against the
classical ICA baseline. The evaluation has two aspects: (i) the effect of
denoising on downstream classification, and (ii) the ability of \saddpm{} to
preserve subject-specific signal structure.

\subsubsection{Downstream Classification}
Tables~\ref{tab:ica} and~\ref{tab:saddpm} report the downstream classification
accuracy (\%) of the cross-subject train--test combinations after ICA and
\saddpm{} denoising, respectively. For each combination, the denoised 2\,s
segments are fed to a downstream four-class motor-imagery classifier, and we
report its classification accuracy; the same classifier and training protocol
are applied to both ICA- and \saddpm{}-denoised data to ensure a fair
comparison.

\begin{table}[t]
  \caption{Downstream classification accuracy (\%) after ICA denoising. Rows are
  the source (training) subject; columns are the target (test) subject. The
  diagonal is the within-subject case; ``Mean'' is the average over source
  subjects for each target.}
  \label{tab:ica}
  \footnotesize
  \setlength{\tabcolsep}{4.2pt}
  \begin{tabular}{c*{9}{c}}
    \toprule
    Source\,$\downarrow$\;/\;Target\,$\rightarrow$ & s1 & s2 & s3 & s4 & s5 & s6 & s7 & s8 & s9 \\
    \midrule
    s1 & 89.29 & 39.29 & 25.93 & 50.00 & 40.74 & 27.27 & 62.96 & 29.63 & 36.00 \\
    s2 & 82.61 & 92.86 & 33.33 & 41.67 & 33.33 & 31.82 & 33.33 & 37.04 & 36.00 \\
    s3 & 53.57 & 85.71 & 81.48 & 58.33 & 33.33 & 54.55 & 44.44 & 33.33 & 48.00 \\
    s4 & 53.57 & 46.43 & 81.48 & 91.67 & 33.33 & 54.55 & 37.04 & 37.04 & 44.00 \\
    s5 & 42.86 & 35.71 & 37.04 & 91.67 & 88.89 & 54.55 & 40.74 & 33.33 & 40.00 \\
    s6 & 50.00 & 39.29 & 37.04 & 50.00 & 77.78 & 90.91 & 44.44 & 33.33 & 44.00 \\
    s7 & 57.14 & 50.00 & 37.04 & 41.67 & 55.56 & 90.91 & 74.07 & 33.33 & 40.00 \\
    s8 & 39.29 & 39.29 & 44.44 & 50.00 & 59.26 & 45.45 & 66.67 & 74.07 & 44.00 \\
    s9 & 32.14 & 32.14 & 37.04 & 45.83 & 44.44 & 45.45 & 33.33 & 74.07 & 84.00 \\
    \midrule
    Mean & 55.61 & 51.19 & 46.09 & 57.87 & 51.85 & 55.05 & 48.56 & 42.80 & 46.22 \\
    \bottomrule
  \end{tabular}
\end{table}

\begin{table}[t]
  \caption{Downstream classification accuracy (\%) after \saddpm{} denoising.
  Layout as in Table~\ref{tab:ica}.}
  \label{tab:saddpm}
  \footnotesize
  \setlength{\tabcolsep}{4.2pt}
  \begin{tabular}{c*{9}{c}}
    \toprule
    Source\,$\downarrow$\;/\;Target\,$\rightarrow$ & s1 & s2 & s3 & s4 & s5 & s6 & s7 & s8 & s9 \\
    \midrule
    s1 & 85.59 & 46.43 & 46.43 & 52.00 & 37.04 & 31.82 & 68.22 & 37.04 & 38.46 \\
    s2 & 85.71 & 90.01 & 34.28 & 44.00 & 44.44 & 37.27 & 32.14 & 44.18 & 35.77 \\
    s3 & 60.71 & 92.86 & 82.11 & 58.00 & 58.15 & 55.00 & 49.29 & 51.85 & 44.62 \\
    s4 & 52.86 & 47.21 & 78.57 & 88.45 & 48.15 & 50.00 & 38.57 & 40.74 & 46.15 \\
    s5 & 45.86 & 35.71 & 37.04 & 91.67 & 85.47 & 52.00 & 42.81 & 36.12 & 46.15 \\
    s6 & 52.14 & 35.00 & 46.43 & 44.44 & 79.12 & 88.00 & 49.15 & 33.33 & 50.00 \\
    s7 & 55.71 & 50.00 & 42.16 & 52.00 & 55.56 & 90.91 & 76.12 & 43.10 & 47.28 \\
    s8 & 39.29 & 39.29 & 44.44 & 47.25 & 59.26 & 45.45 & 70.37 & 78.88 & 44.00 \\
    s9 & 25.15 & 35.71 & 32.14 & 38.20 & 44.44 & 45.45 & 34.07 & 73.08 & 85.15 \\
    \midrule
    Mean & 55.89 & 52.47 & 49.29 & 57.33 & 56.85 & 55.10 & 51.19 & 48.70 & 48.62 \\
    \bottomrule
  \end{tabular}
\end{table}

\begin{table}[t]
  \caption{Aggregate downstream classification summary. ``All'' averages all
  $81$ source--target entries; ``off-diagonal'' excludes the within-subject
  entries and therefore isolates cross-subject transfer; ``diagonal'' averages
  the $9$ within-subject entries.}
  \label{tab:acc_summary}
  \small
  \begin{tabular}{lccc}
    \toprule
    Metric & ICA & \saddpm{} & $\Delta$ \\
    \midrule
    All source--target pairs & 50.58 & 52.83 & $+2.25$ \\
    Off-diagonal cross-subject pairs & 46.25 & 48.88 & $+2.63$ \\
    Diagonal within-subject pairs & 85.25 & 84.42 & $-0.83$ \\
    Std. of per-target means & 4.76 & 3.35 & $-1.41$ \\
    \bottomrule
  \end{tabular}
\end{table}

Overall, \saddpm{} attains more stable cross-subject classification than ICA.
Averaging the per-target means (the ``Mean'' row), \saddpm{} reaches
$52.83\%$ versus $50.58\%$ for ICA, and it improves the mean accuracy for eight
of the nine target subjects, with s4 the only exception and the gain on s6
marginal. As summarized in Table~\ref{tab:acc_summary}, the improvement is
concentrated in off-diagonal source--target pairs: \saddpm{} improves the
cross-subject average from $46.25\%$ to $48.88\%$ ($+2.63$ points), whereas the
within-subject diagonal average decreases slightly from $85.25\%$ to $84.42\%$
($-0.83$ points). This pattern is consistent with the goal of improving
cross-subject robustness rather than optimizing the session-separated
within-subject case.

The improvement is most evident when the source subject has transferable signal
structure. With s3 as the source subject, for instance, \saddpm{} raises
accuracy on several targets substantially (e.g., s5: $33.33\%\rightarrow58.15\%$
and s8: $33.33\%\rightarrow51.85\%$), indicating that its denoised
representations transfer well. By contrast, when s9 is the source subject, both
methods perform poorly across targets, suggesting a higher proportion of
non-additive artifacts or spectral aliasing in s9; in this case neither method
holds a clear advantage.

The absolute gains are modest, but \saddpm{} exhibits smaller fluctuations across
target subjects: the standard deviation of the per-target mean accuracy is
$3.35$ for \saddpm{} versus $4.76$ for ICA. For target subjects with large
distribution gaps (e.g., s8 and s9 as the test subject), the transductively
adapted subject embedding provides a mechanism for adjusting the reverse
denoising trajectory to the target signal distribution. These results suggest
that subject-aware diffusion can stabilize downstream representations under
cross-subject distribution shift. Because BCI-IV-2a does not provide
artifact-free ground truth, however, this phase should be interpreted as
indirect evidence based on downstream decoding rather than a definitive
artifact-level denoising evaluation.

\subsubsection{Subject-Consistency Analysis}
To examine how faithfully \saddpm{}-generated signals reflect real signals, we
compute subject-correlation matrices (Table~\ref{tab:corr}): panel~(a) reports
correlations within the original BCI-IV signals, and panel~(b) reports
correlations between \saddpm{}-generated samples and the real signals. Each entry
is the average correlation between the (per-channel, trial-averaged) signals of
the two corresponding subjects, averaged over channels, so that the diagonal
measures within-subject self-similarity and the off-diagonal entries measure
cross-subject similarity.

\begin{table*}[t]
  \caption{Subject-correlation analysis. (a) Within-subject and cross-subject
  correlation coefficients of the original BCI-IV EEG. (b) Correlation
  coefficients between \saddpm{}-generated samples and the real BCI-IV EEG. In
  both panels, the diagonal (within-subject) coefficients are generally the
  largest in their rows, suggesting retention of subject-specific signal
  structure.}
  \label{tab:corr}
  \footnotesize
  \setlength{\tabcolsep}{5pt}
  \centering
  \par\smallskip
  (a) Original BCI-IV EEG\par\smallskip
  \begin{tabular}{c*{9}{c}}
    \toprule
    Subj.\,$\downarrow$/\,$\rightarrow$ & s1 & s2 & s3 & s4 & s5 & s6 & s7 & s8 & s9 \\
    \midrule
    s1 & 0.102 & 0.047 & 0.055 & 0.059 & 0.045 & 0.051 & 0.050 & 0.042 & 0.043 \\
    s2 & 0.048 & 0.074 & 0.040 & 0.038 & 0.033 & 0.041 & 0.045 & 0.037 & 0.038 \\
    s3 & 0.056 & 0.040 & 0.080 & 0.048 & 0.037 & 0.043 & 0.048 & 0.037 & 0.039 \\
    s4 & 0.059 & 0.038 & 0.048 & 0.088 & 0.041 & 0.042 & 0.054 & 0.043 & 0.044 \\
    s5 & 0.045 & 0.033 & 0.037 & 0.041 & 0.058 & 0.034 & 0.039 & 0.031 & 0.034 \\
    s6 & 0.053 & 0.041 & 0.043 & 0.042 & 0.033 & 0.071 & 0.042 & 0.034 & 0.038 \\
    s7 & 0.050 & 0.046 & 0.048 & 0.041 & 0.039 & 0.043 & 0.092 & 0.042 & 0.050 \\
    s8 & 0.042 & 0.037 & 0.036 & 0.053 & 0.031 & 0.054 & 0.042 & 0.069 & 0.032 \\
    s9 & 0.050 & 0.038 & 0.038 & 0.043 & 0.034 & 0.038 & 0.043 & 0.032 & 0.066 \\
    \bottomrule
  \end{tabular}
  \par\medskip
  (b) \saddpm{}-generated samples vs.\ real BCI-IV EEG\par\smallskip
  \begin{tabular}{c*{9}{c}}
    \toprule
    Gen.\,$\downarrow$/\,Real\,$\rightarrow$ & s1 & s2 & s3 & s4 & s5 & s6 & s7 & s8 & s9 \\
    \midrule
    s1 & 0.109 & 0.071 & 0.077 & 0.080 & 0.066 & 0.073 & 0.072 & 0.063 & 0.070 \\
    s2 & 0.073 & 0.090 & 0.067 & 0.065 & 0.059 & 0.066 & 0.068 & 0.056 & 0.061 \\
    s3 & 0.074 & 0.065 & 0.096 & 0.072 & 0.060 & 0.069 & 0.070 & 0.060 & 0.063 \\
    s4 & 0.077 & 0.063 & 0.073 & 0.100 & 0.063 & 0.068 & 0.074 & 0.064 & 0.064 \\
    s5 & 0.068 & 0.060 & 0.065 & 0.067 & 0.076 & 0.062 & 0.063 & 0.055 & 0.059 \\
    s6 & 0.083 & 0.066 & 0.069 & 0.067 & 0.060 & 0.088 & 0.068 & 0.057 & 0.061 \\
    s7 & 0.072 & 0.069 & 0.079 & 0.075 & 0.071 & 0.069 & 0.103 & 0.066 & 0.066 \\
    s8 & 0.066 & 0.064 & 0.065 & 0.069 & 0.058 & 0.062 & 0.066 & 0.082 & 0.058 \\
    s9 & 0.071 & 0.064 & 0.067 & 0.067 & 0.060 & 0.065 & 0.067 & 0.056 & 0.082 \\
    \bottomrule
  \end{tabular}
\end{table*}

In the \saddpm{}-generated samples, the within-subject correlations are
generally higher than the cross-subject correlations, matching the structure of
the real EEG. For example, subjects s1, s4, and s7 have within-subject
coefficients of $0.109$, $0.100$, and $0.103$, respectively, all larger than
their cross-subject values. This suggests that \saddpm{} retains subject-specific
signal morphology while learning the global EEG distribution. We interpret this
matrix as a subject-consistency check rather than evidence of biometric identity
preservation or complete artifact-level fidelity: correlation alone cannot
establish that all physiological task-relevant components are preserved. The
result nevertheless supports the use of subject-aware conditioning for
personalized EEG preprocessing and downstream decoding.

The \emph{absolute} coefficients are small ($\approx 0.05$--$0.11$), which is
expected and does not weaken this reading. Two preprocessing choices deliberately
remove the dominant, subject-\emph{invariant} sources of correlation: per-window
$z$-scoring discards amplitude and DC scale (a large part of raw inter-channel
correlation), and averaging over trials attenuates activity that is not
phase-locked across trials. What survives is the weak, subject-\emph{specific}
component of the waveform---individual spectral shape and spatial distribution
arising from anatomy and electrode placement (Section~\ref{sec:related_subject})---
which is intrinsically low in magnitude on 2\,s motor-imagery windows and is
further diluted by channel averaging. The physiologically meaningful quantity is
therefore the \emph{relative} structure (each row's diagonal exceeding its
off-diagonal entries), not the absolute level: a subject resembles itself more
than it resembles others in exactly the residual, individual-specific signal that
subject-aware conditioning is meant to model.

\section{Paired Ground-Truth Denoising}
\label{sec:exp_phase2}

The study so far evaluates denoising \emph{indirectly}, through downstream
classification, because BCI-IV-2a provides no artifact-free ground truth. To
measure denoising fidelity \emph{directly}, we additionally adopt the paired
protocol of EEGdenoiseNet~\cite{zhang2021eegdenoisenet}: a clean EEG segment is
combined with a known artifact---ocular (EOG) or myogenic (EMG)---at a controlled
signal-to-noise ratio (SNR), so the recovered signal is scored against the true
clean signal. We report the correlation coefficient with the clean signal
(CC$\uparrow$), averaged over the SNR range $-7$ to $2$\,dB. These experiments use
the conditional variant \mbox{SADDPM-Cond} (Section~\ref{sec:method_cond}); we study a
single-channel benchmark and a multi-channel setting in which the EEGdenoiseNet
artifacts are injected into the 22-channel windows with a physiological spatial
topography (ocular frontal, myogenic lateral), which also provides subject labels
for the subject-ablation study.

\paragraph{Setup.} \mbox{SADDPM-Cond} is trained on synthetic pairs: a clean EEG
segment is combined with an EEGdenoiseNet artifact at an SNR drawn uniformly over
$-7$ to $2$\,dB (a low-SNR regime, $-16$ to $-6$\,dB, is used where noted), and the
network learns to recover the clean segment. The optimizer and schedule follow
Section~\ref{sec:exp_setup} (Adam, $1\times10^{-4}$, cosine, batch~$64$, $T=1000$,
$128$-dimensional subject embedding with $L_2$ weight decay $1\times10^{-4}$). In
the multi-channel setting, artifacts are injected into all $22$ channels with a
physiological spatial topography and the denoiser is trained jointly over
channels; pairs are split by trial into disjoint train and test sets. For the
subject ablation, the ``wrong'' embedding $\mathbf{e}(s')$ is that of a different,
randomly chosen subject, and the ``null'' embedding is the dedicated no-subject
slot; inference uses full conditional generation throughout.

\paragraph{The decisive design choice.} For the conditional denoiser, predicting
the clean signal directly ($\mathbf{x}_0$-prediction) rather than the noise is the
main lever: on the single-channel EOG benchmark it raises CC from $0.848$ to
$0.901$. Combined with full conditional generation
(Section~\ref{sec:method_cond}), this makes \mbox{SADDPM-Cond} competitive with
supervised CNN denoisers on ocular-artifact removal.

\subsection{Single-Channel Benchmark}
On the standard single-channel EEGdenoiseNet benchmark
(Table~\ref{tab:p2_single}, Figure~\ref{fig:f2}), \mbox{SADDPM-Cond} is
\emph{competitive} with supervised CNN denoisers on ocular artifacts (EOG CC
$0.901$ vs.\ SimpleCNN $0.912$ and ComplexCNN $0.909$) and \emph{trails} CNN
baselines more clearly on the harder broadband myogenic artifact (EMG CC
$0.680$ vs.\ SimpleCNN $0.742$ and the best reported baseline, NovelCNN,
$0.831$). We therefore claim viability, not superiority: a diffusion denoiser can
approach standard supervised CNNs on EOG, while EMG denoising remains more
challenging. Example denoised traces
(Figure~\ref{fig:f1}) show the recovered signal closely following the clean
ground truth while the ocular drift is suppressed.

\begin{table}[t]
  \caption{Single-channel denoising on the EEGdenoiseNet benchmark
  (CC$\uparrow$ with the clean signal, averaged over $-7$ to $2$\,dB). CNN
  baselines are from EEGdenoiseNet~\cite{zhang2021eegdenoisenet}.}
  \label{tab:p2_single}
  \small
  \begin{tabular}{lcc}
    \toprule
    Method & EOG CC & EMG CC \\
    \midrule
    Noisy (input) & 0.500 & 0.506 \\
    SimpleCNN & \textbf{0.912} & 0.742 \\
    ComplexCNN & 0.909 & 0.691 \\
    NovelCNN & 0.878 & \textbf{0.831} \\
    \mbox{SADDPM-Cond} (ours) & 0.901 & 0.680 \\
    \bottomrule
  \end{tabular}
\end{table}

\begin{figure}[t]
  \centering
  \includegraphics[width=\linewidth]{figures/F2_metrics_vs_snr_EOG}
  \caption{Single-channel denoising on the EEGdenoiseNet benchmark (EOG).
  RRMSE$_{\text{temporal}}$, RRMSE$_{\text{spectral}}$, and correlation
  coefficient (CC) versus input SNR for \mbox{SADDPM-Cond} and the EEGdenoiseNet CNN
  baselines; ``Noisy'' is the unprocessed input. \mbox{SADDPM-Cond} tracks the
  supervised CNNs across the SNR range, establishing diffusion as a viable EEG
  denoiser.}
  \label{fig:f2}
  \Description{Three line plots versus input SNR from -7 to 2 dB: RRMSE-temporal,
  RRMSE-spectral, and correlation coefficient. The SADDPM-Cond curve closely
  follows the SimpleCNN, ComplexCNN, and NovelCNN curves, all much better than
  the noisy-input reference.}
\end{figure}

\begin{figure}[t]
  \centering
  \includegraphics[width=0.85\linewidth]{figures/F1_waveforms_EOG}
  \caption{Example EOG-contaminated segments at four input SNRs. The
  \mbox{SADDPM-Cond}-denoised signal (blue) recovers the clean ground truth (black)
  and suppresses the ocular drift present in the input (red).}
  \label{fig:f1}
  \Description{Four stacked time-series panels at SNR -7, -5, -3, and 1 dB. In
  each panel the denoised trace overlaps the clean trace, while the contaminated
  trace shows large low-frequency ocular drift, most visible after 1.5 seconds.}
\end{figure}

\subsection{Multi-Channel Joint Denoising}
Denoising all 22 channels jointly lets \mbox{SADDPM-Cond} exploit the artifact's
spatial topography (Proposition~\ref{prop:joint}), reaching CC $\approx 0.99$ for
both artifact types (Table~\ref{tab:p2_multi}, Figure~\ref{fig:f5}), well above
per-channel CNN baselines, which operate on one channel at a time by
construction.

\begin{table}[t]
  \caption{Multi-channel joint denoising (CC$\uparrow$), identical test windows.
  The per-channel SimpleCNN is applied independently to each channel.}
  \label{tab:p2_multi}
  \small
  \begin{tabular}{lcc}
    \toprule
    Method & EOG CC & EMG CC \\
    \midrule
    Noisy (input) & 0.662 & 0.542 \\
    Per-channel SimpleCNN~\cite{zhang2021eegdenoisenet} & 0.844 & 0.853 \\
    \mbox{SADDPM-Cond} (multi-channel, ours) & \textbf{0.988} & \textbf{0.994} \\
    \bottomrule
  \end{tabular}
\end{table}

\begin{figure}[t]
  \centering
  \includegraphics[width=\linewidth]{figures/F5_multichannel_EOG}
  \caption{Multi-channel joint denoising (EOG). Left: the frontal-dominant
  spatial topography of the injected ocular artifact. Right: contaminated,
  \mbox{SADDPM-Cond}-denoised, and clean 22-channel windows (channel $\times$ time);
  the spatially structured artifact is removed and the denoised output matches
  the ground truth.}
  \label{fig:f5}
  \Description{A scalp topomap showing a frontal-dominant ocular artifact,
  followed by three channel-by-time heatmaps: the contaminated window shows a
  strong frontal artifact band, the denoised window removes it and matches the
  clean window.}
\end{figure}

\subsection{Subject-Aware Conditioning Is Functionally Relevant}
\label{sec:exp_subject_loadbearing}
To test whether the subject embedding contributes materially, we denoise each segment
three ways---with the \emph{correct} embedding $\mathbf{e}(s)$, a \emph{wrong}
embedding $\mathbf{e}(s')$, and a \emph{null} embedding---and measure the drop
$\Delta\mathrm{CC}$ between correct and wrong (Table~\ref{tab:p2_subject},
Figure~\ref{fig:f3}). When
the artifact is \emph{shared} across subjects, the embedding has negligible effect
($\Delta\mathrm{CC}\approx 0$): the clean signal is fully recoverable from the
input, so the subject index provides no additional predictive information. When
each subject's artifact occupies a \emph{distinct spatial subspace}
(physiologically plausible, reflecting anatomy and electrode differences), the
embedding becomes functionally relevant---a wrong embedding removes the wrong
subspace and CC collapses from
$0.993$ to $0.828$ ($\Delta\mathrm{CC}=+0.165$), rising to $+0.273$ at low SNR.
The effect is robust across artifact-specificity strength
($\Delta\mathrm{CC}=+0.080/+0.165/+0.139$ at topography gain $0.3/0.6/1.0$) and
holds for EMG ($+0.087$).

\begin{table}[t]
  \caption{Subject ablation for \mbox{SADDPM-Cond} (EOG, CC$\uparrow$): denoising with
  the correct vs.\ a wrong vs.\ a null subject embedding. $\Delta\mathrm{CC}$ is
  correct$-$wrong; $\Delta\mathrm{CC}\ge 0.02$ indicates a functionally relevant embedding.}
  \label{tab:p2_subject}
  \small
  \begin{tabular}{lcccc}
    \toprule
    Regime & correct $\mathbf{e}(s)$ & wrong $\mathbf{e}(s')$ & null & $\Delta\mathrm{CC}$ \\
    \midrule
    Shared artifact & 0.988 & 0.987 & 0.982 & $+0.001$ \\
    Subject-specific artifact & 0.993 & 0.828 & 0.865 & $\mathbf{+0.165}$ \\
    \quad + low SNR & 0.972 & 0.699 & 0.791 & $\mathbf{+0.273}$ \\
    \bottomrule
  \end{tabular}
\end{table}

\begin{figure}[t]
  \centering
  \includegraphics[width=\linewidth]{figures/F3_subject_swap_heatmap}
  \caption{Is the subject embedding functionally relevant? Each cell is the denoising CC
  when subject $i$'s data is denoised with subject $j$'s embedding (diagonal $=$
  correct subject). With a shared artifact (left) the embedding is irrelevant
  (correct $\approx$ wrong); with subject-specific artifacts (right) only the
  correct subject's embedding recovers the signal---a bright diagonal against
  darker off-diagonal entries---showing that subject conditioning
  contributes materially when artifacts are subject-specific (cf.\
  Table~\ref{tab:p2_subject} and Proposition~\ref{prop:obs}).}
  \label{fig:f3}
  \Description{Two heatmaps. Left (shared artifact): nearly uniform high CC
  regardless of which subject embedding is used. Right (subject-specific
  artifact): a bright high-CC diagonal where the correct embedding is used,
  against lower-CC off-diagonal cells where a wrong embedding is used.}
\end{figure}

This is the empirical counterpart of Proposition~\ref{prop:obs}: the value of
subject conditioning equals the between-subject variance of the optimal denoiser,
which is zero for shared artifacts and large for subject-specific ones. We
therefore state the subject-aware claim with its condition: conditioning is
functionally relevant \emph{when artifacts are subject-specific}, which is
physiologically the typical case.

\paragraph{Scope.} These results establish \emph{viability}, not superiority. The
multi-channel and subject-specific settings use semi-synthetic data (artifacts
injected with modeled topographies), as is standard for EEGdenoiseNet-style
paired evaluation; the subject-specific topography is a physiologically-motivated
modeling choice. We do not claim that diffusion is the best learned denoiser, nor
that subject-awareness uniquely requires diffusion; the contribution is that a
subject-aware diffusion denoiser is effective and competitive, and that its
subject mechanism is functionally relevant under subject-specific artifacts.
